%% ============================================================
%% Chapter 1 — Introduction and aim of the thesis
%% ============================================================
\chapter{Introduction and background of the Thesis}\label{chapter:1_Introduction}
The continuous development of technology \dots \: This chapter should present the problem to be analysed or solved within the thesis, together with a justification for undertaking the topic. It should briefly discuss the current state of the art and existing solutions related to the selected issue.

It is recommended to include a theoretical description of the problem, definitions of basic concepts, and a review of existing solutions - highlighting their strengths and weaknesses. It is also important to clearly formulate the thesis topic in line with the approved diploma data in the MyGUT (Moja PG) portal. The entire text should be prepared with editorial consistency, linguistic and terminological correctness, and overall coherence.

\section{State of the art}
For many years, there has been a dynamic increase in the number of publications devoted to \dots \: This section should provide a concise literature review covering 10–20 items. It should include both classic works forming the foundations of the field and the most recent publications from the last five years. The aim is to present, briefly yet coherently, the current state of knowledge and solutions related to the thesis topic, together with their initial analysis and assessment.

Citations in the text should be inserted with the command \verb|\cite{key}|, e.g. \cite{DeanGhemawat2004MapReduce}. Bibliography should be prepared in accordance with the guidelines of the \textbf{ISO 690} standard, using the file \verb|_bibliography.bib| where all sources included in the thesis are defined. The references \cite{HorowitzHill2015AoE} will be collected in the \hyperref[page:Bibliography]{Bibliography}, in accordance with editorial requirements and good thesis-writing practice.

The numbering of bibliography entries is automatic and follows the order of appearance in the text \cite{Kalman1960NewApproach}. If another numbering scheme is required, it can be adjusted via the \texttt{biber} backend, e.g., in the file \verb|_formatting.cls|, although the default setting is recommended.

\section{Aim and Scope of the Thesis}
The subject of this thesis is \dots \: The text should be stated clearly and concisely, including the \textbf{topic} and the precise aim of the thesis. If necessary, a thesis statement (hypothesis) should be formulated. Then, briefly present prior work in the field, the adopted technical assumptions, and provide a short preview of the content of the subsequent chapters.

In the introduction, it is worth outlining the current state of knowledge and technical solutions, identifying unexplored areas and key research and design challenges. It should also be justified that the subsequent parts of the thesis develop the topic in an original and well-founded manner.

In the case of multi-author theses, authorship of individual parts must be clearly assigned. One author should prepare each subsection, and the authors of chapters or subsections should be listed in the appropriate place in the document.
